% 共用表格欄位；由摘要及第一章載入，避免修改主檔。
\ifdefined\proposalsetupdone\else
\gdef\proposalsetupdone{1}
\newcolumntype{P}[1]{>{\raggedright\arraybackslash\hspace{0pt}}p{#1}}
% 書目依使用者指定：中文文獻在前，英文文獻在後。
\nocite{taichung_eat_outside,taichung_gender,tdx_planner,tdx_bus,llm_planning_tools,subjective_attributes,critique_preference,bus_dynamic_replanning,llm_planning,ttg,hypertree_planning,bright,hybrid_retrieval,knowledge_graph_llm_recommendation,restaurant_habib,poi_stisan,bus_delay_uncertainty,google_places_policy}
\fi

\begin{abstractzh}
臺中市政府於 2025 年發布的家庭外食分析指出，2024 年底臺中市餐飲業營利事業家數達 22,612 家。對約會、慶生或朋友聚餐的使用者而言，多樣的餐廳選擇除了需要符合用餐情境，也需要實際方便前往。另一方面，近年的規劃研究已將自然語言理解與可驗證的限制處理結合，提供整合用餐需求與交通分析的技術基礎~\cite{taichung_eat_outside,llm_planning_tools}。

然而，評論與用餐情境在語意上相符，不代表餐廳在交通上適合使用者。若餐廳需要長時間候車、步行或多次轉乘，推薦結果可能仍不實用；若在選定餐廳後才附交通路線，交通條件也無法回饋到餐廳的取捨。因此，本研究關注如何在語意候選形成後，將交通可行性與便利性納入餐廳的最終選擇。

為此，本研究提出一個考量臺中公車交通便利性的情境感知餐廳推薦 AI Agent System，先將使用者情境與餐廳評論語意對齊，再取得候選餐廳的公車搭乘方案，篩選並排序最適合的餐廳。

系統以語言模型整理需求，結合向量或混合檢索召回餐廳，以知識圖譜保存餐廳屬性與來源片段，並透過 TDX 取得路線與可用的公車到站資料。分析引擎計算步行、候車、乘車與轉乘負擔，建立抵達及用餐時間線，再以通過、不通過、不確定三值查核預算、營業與搭乘條件。最終輸出餐廳、適合該情境的評論依據、交通比較與搭乘方案；推薦理由依判定紀錄產生。

本研究目前為提案階段，預期建立一個能說明情境依據與交通取捨的餐廳推薦原型，在保留語意相符餐廳的前提下，結合交通篩選與排序，提供適合用餐情境且方便前往的推薦。

\begin{flushleft}
關鍵詞: \textbf{餐廳推薦、評論語意、情境感知推薦、臺中公車、向量檢索、知識圖譜、三值查核。}
\end{flushleft}
\end{abstractzh}
